\documentclass[a4paper,12pt,twoside]{memoir}

% Castellano
\usepackage[spanish,es-tabla]{babel}
\selectlanguage{spanish}
\usepackage[utf8]{inputenc}
\usepackage[T1]{fontenc}
\usepackage{lmodern} % scalable font
\usepackage{microtype}
\usepackage{placeins}
\usepackage{microtype}
\usepackage{float}
\usepackage{adjustbox}
\usepackage{graphicx}
\usepackage{tabularx}
\usepackage{pdfpages}


\RequirePackage{booktabs}
\RequirePackage[table]{xcolor}
\RequirePackage{xtab}
\RequirePackage{multirow}

% Links
\PassOptionsToPackage{hyphens}{url}\usepackage[colorlinks]{hyperref}
\hypersetup{
	allcolors = {red}
}

% Ecuaciones
\usepackage{amsmath}

% Rutas de fichero / paquete
\newcommand{\ruta}[1]{{\sffamily #1}}

% Párrafos
\nonzeroparskip

% Huérfanas y viudas
\widowpenalty100000
\clubpenalty100000

% Evitar solapes en el header
\nouppercaseheads

% Imagenes
\usepackage{graphicx}
\newcommand{\imagen}[2]{
	\begin{figure}[!h]
		\centering
		\includegraphics[width=0.9\textwidth]{#1}
		\caption{#2}\label{fig:#1}
	\end{figure}
	\FloatBarrier
}

\newcommand{\imagenflotante}[2]{
	\begin{figure}%[!h]
		\centering
		\includegraphics[width=0.9\textwidth]{#1}
		\caption{#2}\label{fig:#1}
	\end{figure}
}



% El comando \figura nos permite insertar figuras comodamente, y utilizando
% siempre el mismo formato. Los parametros son:
% 1 -> Porcentaje del ancho de página que ocupará la figura (de 0 a 1)
% 2 --> Fichero de la imagen
% 3 --> Texto a pie de imagen
% 4 --> Etiqueta (label) para referencias
% 5 --> Opciones que queramos pasarle al \includegraphics
% 6 --> Opciones de posicionamiento a pasarle a \begin{figure}
\newcommand{\figuraConPosicion}[6]{%
  \setlength{\anchoFloat}{#1\textwidth}%
  \addtolength{\anchoFloat}{-4\fboxsep}%
  \setlength{\anchoFigura}{\anchoFloat}%
  \begin{figure}[#6]
    \begin{center}%
      \Ovalbox{%
        \begin{minipage}{\anchoFloat}%
          \begin{center}%
            \includegraphics[width=\anchoFigura,#5]{#2}%
            \caption{#3}%
            \label{#4}%
          \end{center}%
        \end{minipage}
      }%
    \end{center}%
  \end{figure}%
}

%
% Comando para incluir imágenes en formato apaisado (sin marco).
\newcommand{\figuraApaisadaSinMarco}[5]{%
  \begin{figure}%
    \begin{center}%
    \includegraphics[angle=90,height=#1\textheight,#5]{#2}%
    \caption{#3}%
    \label{#4}%
    \end{center}%
  \end{figure}%
}
% Para las tablas
\newcommand{\otoprule}{\midrule [\heavyrulewidth]}
%
% Nuevo comando para tablas pequeñas (menos de una página).
\newcommand{\tablaSmall}[5]{%
 \begin{table}
  \begin{center}
   \rowcolors {2}{gray!35}{}
   \begin{tabular}{#2}
    \toprule
    #4
    \otoprule
    #5
    \bottomrule
   \end{tabular}
   \caption{#1}
   \label{tabla:#3}
  \end{center}
 \end{table}
}

%
%Para el float H de tablaSmallSinColores
\usepackage{float}

%
% Nuevo comando para tablas pequeñas (menos de una página).
\newcommand{\tablaSmallSinColores}[5]{%
 \begin{table}[H]
  \begin{center}
   \begin{tabular}{#2}
    \toprule
    #4
    \otoprule
    #5
    \bottomrule
   \end{tabular}
   \caption{#1}
   \label{tabla:#3}
  \end{center}
 \end{table}
}

\newcommand{\tablaApaisadaSmall}[5]{%
\begin{landscape}
  \begin{table}
   \begin{center}
    \rowcolors {2}{gray!35}{}
    \begin{tabular}{#2}
     \toprule
     #4
     \otoprule
     #5
     \bottomrule
    \end{tabular}
    \caption{#1}
    \label{tabla:#3}
   \end{center}
  \end{table}
\end{landscape}
}

%
% Nuevo comando para tablas grandes con cabecera y filas alternas coloreadas en gris.
\newcommand{\tabla}[6]{%
  \begin{center}
    \tablefirsthead{
      \toprule
      #5
      \otoprule
    }
    \tablehead{
      \multicolumn{#3}{l}{\small\sl continúa desde la página anterior}\\
      \toprule
      #5
      \otoprule
    }
    \tabletail{
      \hline
      \multicolumn{#3}{r}{\small\sl continúa en la página siguiente}\\
    }
    \tablelasttail{
      \hline
    }
    \bottomcaption{#1}
    \rowcolors {2}{gray!35}{}
    \begin{xtabular}{#2}
      #6
      \bottomrule
    \end{xtabular}
    \label{tabla:#4}
  \end{center}
}

%
% Nuevo comando para tablas grandes con cabecera.
\newcommand{\tablaSinColores}[6]{%
  \begin{center}
    \tablefirsthead{
      \toprule
      #5
      \otoprule
    }
    \tablehead{
      \multicolumn{#3}{l}{\small\sl continúa desde la página anterior}\\
      \toprule
      #5
      \otoprule
    }
    \tabletail{
      \hline
      \multicolumn{#3}{r}{\small\sl continúa en la página siguiente}\\
    }
    \tablelasttail{
      \hline
    }
    \bottomcaption{#1}
    \begin{xtabular}{#2}
      #6
      \bottomrule
    \end{xtabular}
    \label{tabla:#4}
  \end{center}
}

%
% Nuevo comando para tablas grandes sin cabecera.
\newcommand{\tablaSinCabecera}[5]{%
  \begin{center}
    \tablefirsthead{
      \toprule
    }
    \tablehead{
      \multicolumn{#3}{l}{\small\sl continúa desde la página anterior}\\
      \hline
    }
    \tabletail{
      \hline
      \multicolumn{#3}{r}{\small\sl continúa en la página siguiente}\\
    }
    \tablelasttail{
      \hline
    }
    \bottomcaption{#1}
  \begin{xtabular}{#2}
    #5
   \bottomrule
  \end{xtabular}
  \label{tabla:#4}
  \end{center}
}



\definecolor{cgoLight}{HTML}{EEEEEE}
\definecolor{cgoExtralight}{HTML}{FFFFFF}

%
% Nuevo comando para tablas grandes sin cabecera.
\newcommand{\tablaSinCabeceraConBandas}[5]{%
  \begin{center}
    \tablefirsthead{
      \toprule
    }
    \tablehead{
      \multicolumn{#3}{l}{\small\sl continúa desde la página anterior}\\
      \hline
    }
    \tabletail{
      \hline
      \multicolumn{#3}{r}{\small\sl continúa en la página siguiente}\\
    }
    \tablelasttail{
      \hline
    }
    \bottomcaption{#1}
    \rowcolors[]{1}{cgoExtralight}{cgoLight}

  \begin{xtabular}{#2}
    #5
   \bottomrule
  \end{xtabular}
  \label{tabla:#4}
  \end{center}
}




\graphicspath{ {./img/} }

% Capítulos
\chapterstyle{bianchi}
\newcommand{\capitulo}[2]{
	\setcounter{chapter}{#1}
	\setcounter{section}{0}
	\setcounter{figure}{0}
	\setcounter{table}{0}
	\chapter*{#2}
	\addcontentsline{toc}{chapter}{#2}
	\markboth{#2}{#2}
}

% Apéndices
\renewcommand{\appendixname}{Apéndice}
\renewcommand*\cftappendixname{\appendixname}

\newcommand{\apendice}[1]{
	%\renewcommand{\thechapter}{A}
	\chapter{#1}
}

\renewcommand*\cftappendixname{\appendixname\ }

% Formato de portada
\makeatletter
\usepackage{xcolor}
\newcommand{\tutor}[1]{\def\@tutor{#1}}
\newcommand{\course}[1]{\def\@course{#1}}
\definecolor{cpardoBox}{HTML}{E6E6FF}
\def\maketitle{
  \null
  \thispagestyle{empty}
  % Cabecera ----------------
\noindent\includegraphics[width=\textwidth]{cabecera}\vspace{1cm}%
  \vfill
  % Título proyecto y escudo informática ----------------
  \colorbox{cpardoBox}{%
    \begin{minipage}{.8\textwidth}
      \vspace{.5cm}\Large
      \begin{center}
      \textbf{TFG del Grado en Ingeniería Informática}\vspace{.6cm}\\
      \textbf{\LARGE\@title{}}
      \end{center}
      \vspace{.2cm}
    \end{minipage}

  }%
  \hfill\begin{minipage}{.20\textwidth}
    \includegraphics[width=\textwidth]{escudoInfor}
  \end{minipage}
  \vfill
  % Datos de alumno, curso y tutores ------------------
  \begin{center}%
  {%
    \noindent\LARGE
    Presentado por \@author{}\\ 
    en Universidad de Burgos --- \@date{}\\
    Tutor: \@tutor{}\\
  }%
  \end{center}%
  \null
  \cleardoublepage
  }
\makeatother


% Datos de portada
\title{Implementación de una
	Prueba de Concepto (PoC) de
	Manage Engine PAM para la
	Gestión Eficiente de Acceso
	Privilegiado en Organizaciones \\Documentación Técnica}
\author{Alejandro Villar Solla}
\tutor{Luis Antonio Antolín Sánchez}
\date{\today}

\begin{document}

\maketitle



\cleardoublepage



%%%%%%%%%%%%%%%%%%%%%%%%%%%%%%%%%%%%%%%%%%%%%%%%%%%%%%%%%%%%%%%%%%%%%%%%%%%%%%%%%%%%%%%%



\frontmatter


\clearpage

% Indices
\tableofcontents

\clearpage

\listoffigures

\clearpage

\listoftables

\clearpage

\mainmatter

\appendix

\apendice{Plan de Proyecto Software}

\section{Introducción}

Este documento detalla el Plan de Proyecto para la implementación de PAM360 en NTT Data, cubriendo todos los aspectos desde la planificación inicial hasta la finalización del proyecto. El objetivo es proporcionar una guía clara y detallada que asegure el éxito en la implementación y configuración de la herramienta PAM360 para la gestión de accesos privilegiados.

\section{Objetivos}

El objetivo principal de este proyecto es implementar y configurar PAM360 para mejorar la gestión de accesos privilegiados en NTT Data. Los objetivos específicos incluyen:

\begin{itemize}
	\item Mejorar la seguridad mediante el control y monitoreo de accesos privilegiados.
	\item Optimizar la eficiencia en la gestión de usuarios y contraseñas.
	\item Asegurar la transparencia y trazabilidad de las actividades de los usuarios.
	\item Facilitar la auditoría y el cumplimiento normativo.
	\item Integrar PAM360 con las infraestructuras existentes de NTT Data.
\end{itemize}

\section{Alcance del Proyecto}

El alcance del proyecto abarca la instalación, configuración e integración de PAM360 con el Active Directory de NTT Data, incluyendo:

\begin{itemize}
	\item Revisión de documentación técnica.
	\item Preparación del entorno de pruebas.
	\item Despliegue y configuración de PAM360.
	\item Integración con AD y pruebas funcionales.
	\item Configuración de políticas de seguridad y contraseñas.
	\item Generación de reportes y auditorías.
	\item Formación y capacitación para los administradores del sistema.
\end{itemize}

\section{Cronograma}

El proyecto se desarrollará en las siguientes fases:

\begin{itemize}
	\item \textbf{Fase 1: Revisión de Documentación} (1 semana)
	\item \textbf{Fase 2: Preparación del Entorno} (1 semanas)
	\item \textbf{Fase 3: Implementación} (1 semanas)
	\item \textbf{Fase 4: Configuración y Pruebas Funcionales} (1 semanas)
	\item \textbf{Fase 5: Evaluación y Análisis} (1 semanas)
	\item \textbf{Fase 6: Documentación Final y Capacitación} (2 semana)
\end{itemize}


\section{Recursos}

\subsection{Recursos Humanos}

\begin{itemize}
	\item \textbf{Gestor de Proyecto}: Responsable de la coordinación y seguimiento del proyecto.
	\item \textbf{Especialista en Seguridad}: Encargado de la configuración y pruebas de seguridad.
	\item \textbf{Administrador de Sistemas}: Responsable de la integración con AD y gestión de usuarios.
\end{itemize}

\subsection{Recursos Materiales y de Software}

\begin{itemize}
	\item Servidores virtuales para pruebas y despliegue.
	\item Licencias de software para PAM360 y herramientas de virtualización (vSphere, MRemote).
	\item Documentación técnica y manuales de usuario.
	\item Equipos de red y hardware de soporte.
\end{itemize}


\section{Riesgos y Mitigación}

\begin{table}[H]
	\centering
	\begin{adjustbox}{max width=\textwidth}
		\begin{tabular}{|p{0.5\linewidth}|p{0.5\linewidth}|}
			\hline
			\textbf{Riesgo} & \textbf{Mitigación} \\ \hline
			Problemas de Integración con AD & Realizar pruebas preliminares en un entorno controlado. \\ \hline
			Fallos en la Configuración de PAM360 & Contar con el soporte técnico del proveedor y documentación detallada. \\ \hline
			Resistencia al Cambio por parte del Personal & Implementar programas de formación y capacitación. \\ \hline
			Vulnerabilidades de Seguridad & Actualizar regularmente el software y realizar auditorías de seguridad. \\ \hline
			Problemas de Conectividad & Buena configuración en el firewall de permisos y mantenimiento pertinente. \\ \hline
		\end{tabular}
	\end{adjustbox}
	\caption{Riesgos y Mitigación}
	\label{tab:riesgos_mitigacion}
\end{table}

\section{Control y Seguimiento}

Para asegurar el cumplimiento de los objetivos del proyecto, se llevarán a cabo las siguientes actividades de control y seguimiento:

\begin{itemize}
	\item Reuniones semanales de seguimiento con el equipo de proyecto.
	\item Reportes de avance semanales.
	\item Evaluaciones periódicas de cumplimiento de hitos.
	\item Auditorías internas para verificar la correcta implementación y configuración de PAM360.
\end{itemize}


\section{Conclusiones}

El Plan de Proyecto Software para la implementación de PAM360 en NTT Data establece una guía clara para llevar a cabo todas las fases del proyecto, asegurando una gestión eficiente y segura de los accesos privilegiados en la organización. Con una planificación detallada, asignación adecuada de recursos y mitigación de riesgos, se espera alcanzar todos los objetivos propuestos de manera exitosa. Además, la integración con las infraestructuras existentes y la capacitación del personal son factores clave para el éxito del proyecto y su sostenibilidad a largo plazo.



\apendice{Plan de Calidad del Proyecto}

\section{Introducción}

El Plan de Calidad del Proyecto establece los procedimientos y estándares necesarios para asegurar que la implementación de PAM360 en NTT Data cumpla con los requisitos de calidad definidos. Este documento detalla las actividades de control de calidad que se llevarán a cabo durante el proyecto.

\section{Casos de uso}
A continuación, se describen los principales casos de uso identificados para la implementación de PAM360:

\begin{table}[p]
	\centering
	\begin{tabularx}{\linewidth}{ p{0.21\columnwidth} p{0.71\columnwidth} }
		\toprule
		\textbf{CU-1}    & \textbf{Instalación y Configuración de PAM360}\\
		\toprule
		\textbf{Versión}              & 1.0    \\
		\textbf{Autor}                & Alejandro Villar Solla \\
		\textbf{Requisitos asociados} & RF-01, RF-02 \\
		\textbf{Descripción}          & Este caso de uso describe los pasos para instalar y configurar PAM360 en el entorno virtualizado de NTT Data. \\
		\textbf{Precondición}         & 
		\begin{itemize}
			\item Entorno virtualizado preparado.
			\item Archivo de instalación de PAM360 disponible.
		\end{itemize}\\
		\textbf{Acciones}             &
		\begin{enumerate}
			\item Descargar el archivo de instalación de PAM360 desde el sitio web del proveedor.
			\item Copiar el archivo .exe a la máquina virtual utilizando MRemote.
			\item Ejecutar el instalador y seguir las instrucciones de instalación.
			\item Configurar los parámetros iniciales (puertos, credenciales administrativas).
		\end{enumerate}\\
		\textbf{Postcondición}        & 
		\begin{itemize}
			\item PAM360 instalado y corriendo en localhost:8282.
			\item Acceso a la interfaz web de PAM360.
		\end{itemize}\\
		\textbf{Excepciones}          & 
		\begin{itemize}
			\item Error durante la instalación debido a incompatibilidad de SO o errores comunes.
			\item Falta de permisos administrativos para ejecutar el instalador.
		\end{itemize}\\
		\textbf{Importancia}          & Alta \\
		\bottomrule
	\end{tabularx}
	\caption{CU-1 Instalación y Configuración de PAM360.}
\end{table}

\begin{table}[p]
	\centering
	\begin{tabularx}{\linewidth}{ p{0.21\columnwidth} p{0.71\columnwidth} }
		\toprule
		\textbf{CU-2}    & \textbf{Integración de PAM360 con Active Directory}\\
		\toprule
		\textbf{Versión}              & 1.0    \\
		\textbf{Autor}                & Alejandro Villar Solla \\
		\textbf{Requisitos asociados} & RF-03, RF-04 \\
		\textbf{Descripción}          & Este caso de uso describe el proceso de integración de PAM360 con el Active Directory (AD) de NTT Data para sincronización de usuarios y roles. \\
		\textbf{Precondición}         & 
		\begin{itemize}
			\item PAM360 instalado y configurado.
			\item Acceso administrativo a Active Directory.
		\end{itemize}\\
		\textbf{Acciones}             &
		\begin{enumerate}
			\item Acceder a la interfaz web de PAM360.
			\item Navegar a la configuración de administración y seleccionar "Import From Active Directory".
			\item Ingresar el nombre del dominio AD y las credenciales de la cuenta asociada.
			\item Seleccionar los usuarios a sincronizar.
			\item Confirmar y ejecutar la integración.
		\end{enumerate}\\
		\textbf{Postcondición}        & 
		\begin{itemize}
			\item Usuarios y roles del AD sincronizados en PAM360.
			\item Accesos privilegiados gestionados centralmente desde PAM360.
		\end{itemize}\\
		\textbf{Excepciones}          & 
		\begin{itemize}
			\item Error de autenticación con AD.
		\end{itemize}\\
		\textbf{Importancia}          & Alta \\
		\bottomrule
	\end{tabularx}
	\caption{CU-2 Integración de PAM360 con Active Directory.}
\end{table}

\begin{table}[p]
	\centering
	\begin{tabularx}{\linewidth}{ p{0.21\columnwidth} p{0.71\columnwidth} }
		\toprule
		\textbf{CU-3}    & \textbf{Configuración de Políticas de Seguridad en PAM360}\\
		\toprule
		\textbf{Versión}              & 1.0    \\
		\textbf{Autor}                & Alejandro Villar Solla \\
		\textbf{Requisitos asociados} & RF-05, RF-06 \\
		\textbf{Descripción}          & Este caso de uso describe el procedimiento para configurar las políticas de seguridad dentro de PAM360, incluyendo autenticación y gestión de contraseñas. \\
		\textbf{Precondición}         & 
		\begin{itemize}
			\item PAM360 instalado y configurado.
			\item Acceso administrativo a la consola de PAM360.
		\end{itemize}\\
		\textbf{Acciones}             &
		\begin{enumerate}
			\item Iniciar sesión en la consola web de PAM360.
			\item Navegar a la sección de configuración de administración.
			\item Seleccionar "Políticas de Seguridad" y luego "Gestión de Contraseñas".
			\item Configurar las reglas de complejidad de contraseñas, frecuencia de cambio y políticas de bloqueo de cuentas.
			\item Guardar y aplicar las políticas configuradas.
		\end{enumerate}\\
		\textbf{Postcondición}        & 
		\begin{itemize}
			\item Políticas de seguridad y contraseñas aplicadas y activas en PAM360.
			\item Refuerzo de la seguridad en la gestión de accesos privilegiados.
		\end{itemize}\\
		\textbf{Excepciones}          & 
		\begin{itemize}
			\item Error al guardar las políticas debido a conflictos con configuraciones existentes.
			\item Usuarios no cumplen con las nuevas políticas de contraseña y necesitan soporte.
		\end{itemize}\\
		\textbf{Importancia}          & Alta \\
		\bottomrule
	\end{tabularx}
	\caption{CU-3 Configuración de Políticas de Seguridad en PAM360.}
\end{table}

\section{Objetivos de Calidad}

Los objetivos de calidad del proyecto incluyen:

\begin{itemize}
	\item Garantizar que PAM360 se instale y configure correctamente según las especificaciones del proveedor y los requisitos de NTT Data.
	\item Asegurar que todas las funcionalidades de PAM360 se integren y operen sin problemas con el Active Directory de NTT Data.
	\item Cumplir con los estándares de seguridad y buenas prácticas en la gestión de accesos privilegiados.
	\item Obtener retroalimentación y aprobación de los usuarios finales y administradores del sistema.
\end{itemize}

\section{Estrategia de Calidad}

Para alcanzar los objetivos de calidad, se adoptarán las siguientes estrategias:

\begin{itemize}
	\item \textbf{Revisión y Validación de Requisitos}: Verificar que todos los requisitos del proyecto estén claros, completos y sean factibles.
	\item \textbf{Pruebas Exhaustivas}: Realizar pruebas de instalación, configuración, integración y funcionalidad para asegurar que el sistema opere correctamente.
	\item \textbf{Auditorías de Seguridad}: Ejecutar auditorías para evaluar la seguridad del sistema y la conformidad con las políticas de seguridad de NTT Data.
	\item \textbf{Capacitación de Usuarios}: Asegurar que los usuarios finales y administradores reciban formación adecuada para utilizar y gestionar PAM360 eficientemente.
\end{itemize}

\section{Procesos para Garantizar la Calidad}

\subsection{Revisión de Documentación}

\begin{itemize}
	\item \textbf{Documentación Técnica}: Revisión de manuales, guías de instalación y configuración proporcionados por ManageEngine.
\end{itemize}

\subsection{Pruebas y Validaciones}

\begin{itemize}
	\item \textbf{Pruebas de Instalación}: Verificación de la correcta instalación de PAM360 en el entorno virtualizado.
	\item \textbf{Pruebas de Configuración}: Asegurar que todas las configuraciones necesarias (autenticación, políticas de seguridad, integración con AD) se realicen correctamente.
	\item \textbf{Pruebas Funcionales}: Validación de todas las funcionalidades de PAM360, incluyendo la gestión de usuarios, contraseñas y auditorías.
	\item \textbf{Pruebas de Seguridad}: Realizar pruebas de penetración y vulnerabilidad para identificar y mitigar posibles riesgos de seguridad.
\end{itemize}

\subsection{Control de Cambios}

\begin{itemize}
	\item \textbf{Gestión de Cambios}: Implementar un sistema de control de cambios para manejar las modificaciones en los requisitos, el diseño o la configuración de PAM360.
	\item \textbf{Aprobación de Cambios}: Todos los cambios deben ser revisados y aprobados por el equipo de proyecto antes de su implementación.
\end{itemize}

\section{Criterios de Aceptación}

Los criterios para aceptar el proyecto incluyen:

\begin{itemize}
	\item \textbf{Cumplimiento de Requisitos}: Todos los requisitos del proyecto deben ser cumplidos satisfactoriamente.
	\item \textbf{Validación de Funcionalidades}: Todas las funcionalidades de PAM360 deben operar sin errores y conforme a las especificaciones.
	\item \textbf{Resultados de Pruebas y Auditorías}: Los resultados de las pruebas y auditorías deben demostrar que el sistema es seguro y funcional.
\end{itemize}

\section{Métricas de Calidad}

Para medir la calidad del proyecto, se utilizarán las siguientes métricas:

\begin{itemize}
	\item \textbf{Tasa de Éxito de Pruebas}: Porcentaje de pruebas que se completan exitosamente sin errores.
	\item \textbf{Número de Incidentes de Seguridad}: Cantidad de incidentes de seguridad detectados y corregidos durante el proyecto.
	\item \textbf{Satisfacción de Usuarios}: Nivel de satisfacción de los usuarios finales y administradores con el sistema implementado.
	\item \textbf{Tiempo de Respuesta a Incidentes}: Tiempo promedio para detectar y resolver incidentes o problemas en el sistema.
\end{itemize}

\section{Responsabilidades}

\begin{itemize}
	\item \textbf{Gestor de Calidad}: Responsable de la planificación, ejecución y seguimiento de las actividades de aseguramiento y control de calidad.
	\item \textbf{Equipo de Proyecto}: Participar en las pruebas, auditorías y revisiones de calidad.
\end{itemize}


\include{./tex/C_especificacion_de_permisos}






\end{document}
